%% 09_chopper.tex -- Chapter 9: Chopper Stabilisation and Offset Cancellation
%% Status: OUTLINE (build stage 3). Replace the outline with the full chapter when it is written.
\chapter{Chopper Stabilisation and Offset Cancellation}\label{ch:chopper}
\chapterlead{Chopping moves the signal above the flicker-noise corner, amplifies it there and brings it back, leaving the offset and $1/f$ noise behind.}

\begin{specnote}[Chapter status: outline]
This chapter is an outline. It will be written in build stage 3 (Part II with Labs 3--5). The topics, the design question and the labs below are fixed; the derivations, worked examples and figures are still to come.
\end{specnote}

\begin{objectives}
\begin{itemize}
  \item explain chopping as modulation and demodulation in the frequency domain;
  \item predict the residual offset and the output ripple of a chopper amplifier;
  \item design a ripple-reduction loop;
  \item compare chopping, auto-zeroing and correlated double sampling;
  \item quantify the input-impedance cost of an input chopper.
\end{itemize}
\end{objectives}

\begin{insight}[Design question]
How much offset, ripple and input-impedance loss does chopping leave behind, and how do you suppress each?
\end{insight}

\section{The modulation view of chopping}\label{sec:ch9-1}
Multiplying by a square wave at $f_\mathrm{ch}$ moves the signal to odd harmonics of $f_\mathrm{ch}$; the amplifier's offset and $1/f$ noise stay at baseband and are modulated away by the output chopper.

\section{Residual offset and ripple}\label{sec:ch9-2}
Charge injection and clock feedthrough of the input switches create a residual offset; the up-modulated amplifier offset appears as a triangular ripple at the output.

\section{Ripple-reduction loops}\label{sec:ch9-3}
Sensing the ripple after the output chopper and feeding a correction back; the 60\,nV/$\sqrt{\mathrm{Hz}}$, \qty{60}{\micro\watt} front-end of Yazicioglu et al.~\cite{yazicioglu2007} as the reference design.

\section{Auto-zero and correlated double sampling}\label{sec:ch9-4}
Sampling the offset and subtracting it; the noise folding that makes auto-zero noisier than chopping.

\section{Chopper spikes and input impedance}\label{sec:ch9-5}
An input chopper in front of the input capacitors draws charge every cycle, so the input impedance falls to roughly $1/(2 f_\mathrm{ch} C_\mathrm{in})$; why this matters for dry electrodes and how impedance boosting (\cref{ch:dsl}) recovers it.

\begin{tryit}[Labs for this chapter]
Lab~5 (\cref{sec:lab5}) adds input and output choppers and a DC servo loop to your LNA: $1/f$ removal, ripple, input-impedance loss and \qty{\pm 300}{\milli\volt} offset tolerance.
\end{tryit}

\paragraph{Read with this chapter} Yazicioglu et al.~\cite{yazicioglu2007}, Yazicioglu, Van Hoof and Puers~\cite{yazicioglubook}.
